\subsection{An order-six Morita--grading obstruction}
\label{sec:morita-order6}

The dimension reconstruction above can force a dual category whose
based ring has a short exhaustive parametrization. This supplies two
further exclusions without assuming a pivotal structure on the original
category.

\begin{theorem}\label{thm:morita-order6}
Let $\mathcal C$ be a fusion category, and let $g$ be an invertible
involution fixing a simple object. Set $A=\one\oplus g$. Suppose the
simple Frobenius--Perron dimensions of
$\mathcal D={}_A\mathcal C_A$ are
\[
  \underbrace{1,\ldots,1}_{6},\qquad
  \underbrace{d,\ldots,d}_{6},\qquad d=2+\sqrt5.
\]
Then $\mathcal C$ has a faithful $C_3$ grading. In particular, a fusion
ring with trivial universal grading and the stated forced dual dimension
profile cannot be categorified. This excludes
$\RR{5}{100}$ and $\RR{5}{408}$ over $\CC$.
\end{theorem}

\begin{proof}
The algebra $A$ exists by the fixed-object argument of
Theorem~\ref{thm:involution}. We first exhaust the possible based rings
of $\mathcal D$. Its invertible objects form a group $G$ of order six.
The multiplicity of an invertible $h$ in $\rho\rho^*$, for a
noninvertible simple $\rho$, is one precisely when $h\rho\cong\rho$.
The identity $d^2=1+4d$ and irrationality of $d$ therefore imply that
the stabilizer of $\rho$ is trivial. Thus the noninvertibles are exactly
$h\rho$, $h\in G$. The right action is free for the same reason, so
\begin{equation}\label{eq:morita-six-relations}
 \rho h=\phi(h)\rho,\qquad
 \rho^2=u+\sum_{h\in G}a_hh\rho,\qquad
 \sum_{h\in G}a_h=4,
\end{equation}
where $\phi\in\operatorname{Aut}(G)$, $u\in G$, and
$a_h\in\ZZ_{\ge0}$. Associativity gives
$\phi^2=\operatorname{Ad}(u)$ and $\phi(u)=u$; comparison of
$(\rho^2)h$ with $\rho(\rho h)$ also makes the coefficient function
invariant under
\[
 k\longmapsto\phi^2(h)k\phi(h)^{-1}.
\]
Only $C_6$ and $S_3$ can occur as $G$.

\smallskip\noindent
\emph{The cyclic case.}
The automorphism $\phi$ is either the identity or inversion. Inversion
would make $(a_h)$ constant on cosets of the square subgroup of order
three, contradicting $\sum a_h=4$. Hence $\phi$ is the identity.
Use additive labels $0,\ldots,5$ for $C_6$. Duality gives
$\rho^*=(-u)\rho$ and
\begin{equation}\label{eq:morita-six-domain}
 u\in\ZZ/6\ZZ,\qquad a_k\ge0,\qquad
 \sum_{k=0}^5a_k=4,\qquad a_k=a_{u-k}.
\end{equation}
There are exactly $60$ pairs $(u,a)$ satisfying these conditions.
Each supplies the entire candidate tensor by
\[
 h\cdot k=h+k,\qquad h\cdot(k\rho)=(h+k)\rho,\qquad
 (h\rho)(k\rho)=h+k+u+\sum_j a_j(h+k+j)\rho.
\]
An exact exhaustive calculation partitions the $60$ pairs as follows:
\begin{center}
\begin{tabular}{@{}rl@{}}
\toprule
Number & Exact consequence \\
\midrule
$42$ & A formal-codegree norm obstruction \\
$6$ & A closed based subring isomorphic to $\RR{2}{20}$ \\
$12$ & A faithful $C_3$ grading with neutral invertible involution \\
\bottomrule
\end{tabular}
\end{center}
The first row is a general obstruction here: by
\cite[Lemma~3.1]{SchTwo}, a category with these two distinct
Frobenius--Perron dimensions is Galois conjugate to a pseudo-unitary
category with the same based ring. In that category $\FPdim(\mathcal
D)/f$ must be an algebraic integer for every formal codegree $f$.
The cyclic candidate ring is commutative, so the Casimir characteristic
roots used in the check are its formal codegrees, with no representation
dimension correction. Its global dimension is
\[
 D=6+6d^2=60+24\sqrt5,
 \qquad p_D(t)=t^2-120t+720.
\]
For each of the $42$ cases, the certificate gives the exact minimal
polynomial $q$ of a formal codegree and a prime $p$ for which
\begin{equation}\label{eq:morita-normalized-norm}
 \frac{v_p(|p_D(0)|)}{\deg p_D}
 <\frac{v_p(|q(0)|)}{\deg q}.
\end{equation}
This contradicts integrality of $D/f$: take the norm in any common
number field containing $D$ and $f$. The second row uses the previously
established general exclusion for $\RR{2}{20}$, whose original exact
certificate is replayed, not just its allocated status.
For each of the remaining $12$ pairs the grading is explicitly
\begin{equation}\label{eq:morita-six-grading}
 \deg(h)=h\pmod3,\qquad \deg(h\rho)=h+2u\pmod3.
\end{equation}
Every nonzero structure constant respects this degree assignment.
The unique nonidentity invertible involution, namely $3\in C_6$, has
degree zero.

\smallskip\noindent
\emph{The nonabelian case.}
Every automorphism of $S_3$ is inner. Replacing $\rho$ by an invertible
multiple makes $\phi$ the identity. Then $u$ is central, so $u=1$,
and $a$ is constant on conjugacy classes. Its values
$(a_e,a_t,a_c)$ on the identity, transpositions and three-cycles satisfy
\[
 a_e+3a_t+2a_c=4.
\]
The complete domain and its obstructions are
\begin{center}
\begin{tabular}{@{}ll@{}}
\toprule
$(a_e,a_t,a_c)$ & Obstruction \\
\midrule
$(4,0,0)$ & The rank-two subring $\rho^2=1+4\rho$ \\
$(2,0,1)$ & The closed subring $\RR{2}{20}$ \\
$(0,0,2)$ & The closed subring $\RR{2}{21}$ \\
$(1,1,0)$ & A one-dimensional formal-codegree norm obstruction \\
\bottomrule
\end{tabular}
\end{center}
The first row is excluded by \cite{Ostrik2}. In the next two rows the
subring is generated by $C_3$ and $\rho$; its based isomorphism and the
previously established general certificate are checked exactly.
For the last row, extend the sign representation of $S_3$ by assigning
$\rho$ a root $\lambda$ of $\lambda^2+2\lambda-1=0$. Equation
\eqref{eq:morita-six-relations} verifies this as a one-dimensional
representation of the full ring. Its genuine formal codegree is
\[
 f=\sum_{h\in S_3}1+\sum_{h\in S_3}\lambda^2
   =6(1+\lambda^2)=24\pm12\sqrt2,
 \qquad q(t)=t^2-48t+288.
\]
Here $\rho^*=\rho$. Since $v_2(720)=4$ and $v_2(288)=5$, condition
\eqref{eq:morita-normalized-norm} holds. The two-dimension reduction
again makes this a general categorical obstruction. In particular, this
step does not identify an unweighted Casimir eigenvalue from a
higher-dimensional irreducible representation with a formal codegree.
All nonabelian possibilities have now been excluded.

\smallskip\noindent
\emph{Returning through Morita duality.}
Consequently $\mathcal D$ has one of the faithful gradings
\eqref{eq:morita-six-grading}. Regard $A$ as the distinguished simple
right module in $\mathcal M=\mathcal C_A$, and let
$B=\underline{\operatorname{End}}_{\mathcal D}(A)$. Double Morita
duality identifies $\mathcal C$ with ${}_B\mathcal D_B$, up to the
opposite tensor convention; see \cite[Chapter~7]{EGNO}.
For a simple $A$-bimodule $Y$, its multiplicity in $B$ is
\[
 \dim\operatorname{Hom}_{\mathcal M}(A\otimes_A Y,A)
 =\dim\operatorname{Hom}_{\mathcal M}(Y,A).
\]
There are exactly two simple bimodules restricting to the right module
$A$: the regular bimodule and its twist by the nontrivial character of
$C_2$. Each has multiplicity one. This can also be read directly from
the two unit columns at $A$ in the Gram matrix $I+M_g$ of
Proposition~\ref{prop:dual-module-gram}; off-diagonal entries in that
row vanish. All other simple bimodules have zero multiplicity. Thus
$B=\one\oplus b$, where $b$ is a nontrivial invertible involution.
In the cyclic case $b$ is necessarily $3\in C_6$, hence $B$ lies in
the neutral component.

A connected separable algebra in the neutral component of a faithful
group grading transfers that grading to its bimodule category. Indeed,
the left and right algebra actions preserve each homogeneous support,
and tensor product over the algebra multiplies the degrees. Every
component is nonzero: for a nonzero homogeneous object $Y$, the free
bimodule $B\otimes Y\otimes B$ is nonzero and has the same degree.
It follows that $\mathcal C$ is faithfully $C_3$-graded.

Finally, for $\RR{5}{100}$ the involution is basis label $4$, the
fixed labels are $3,5$, and the exhaustive anti-invariant generator
search has $16$ tuples. For $\RR{5}{408}$ these data are respectively
$1$, $4,5$, and $8$. Each search has exactly one admissible
representation, in which the involution acts as $+I$ on the two
anti-invariant coordinates. The Gram reconstruction forces exactly the
dimension profile in the statement. In both original tensors the
multiplicative closure of the supports of all $XX^*$ is the whole
basis. Thus the adjoint subcategory would be all of $\mathcal C$,
which rules out any nontrivial faithful grading. This proves the two
claimed exclusions.
\end{proof}

\paragraph{Exact finite certificate.}
The implementation is \filename{software/new_morita_order6.py}, with
audit output in
\filename{research/model_continuation/morita_order6_certificates.json}.
The input certificate stores the involution, fixed labels, generator
search size, every admissible anti-invariant representation, and the
reconstructed dual dimensions in exact field coordinates. Verification
reconstructs the exhaustive representation domain, rather than trusting
a proposed list of branches. The finite dual routine enumerates every
composition of $4$ into six nonnegative parts, applies exactly
\eqref{eq:morita-six-domain}, and checks that there are $60$ cases.
It constructs each full tensor, factors its Casimir polynomial exactly,
and checks the recorded norm witness, the full grading identities, or
the closed-subring support and unit-preserving based isomorphism.
The four $S_3$ class triples are independently exhausted. The inherited
certificates for $\RR{2}{20}$ and $\RR{2}{21}$ are replayed before they
are used. Thus the audit output is a reproducible record of a complete
finite necessary-condition calculation, with no unresolved branch
silently discarded.

A separate replay, \filename{research/model_continuation/replay_morita_order6.py},
checks equality with the saved finite audit and rejects mutations of the
involution, branch list, generator count, representation matrices, dual
dimension counts, and input basis labels. Its checks use explicit
exceptions and remain active under \texttt{python -O}.
